% =============================================================================
% CVE-2026-73570 — Xpectra Security Research
% "CVE-2026-73570: Full Lifecycle Evidence, Defender Detections, and
%  Remediation for Zimbra's SNMP-Sink Command Injection"
% Build: pdflatex main.tex (twice)
% =============================================================================
\documentclass[journal]{IEEEtran}
\usepackage[utf8]{inputenc}
\usepackage{graphicx}
\usepackage{booktabs}
\usepackage{xcolor}
\usepackage{listings}
\usepackage{hyperref}
\usepackage{microtype}
\usepackage{subcaption}

\hypersetup{colorlinks=true,linkcolor=blue,citecolor=blue,urlcolor=blue}

\definecolor{codebg}{RGB}{245,247,250}
\definecolor{codekw}{RGB}{0,90,150}
\definecolor{codecm}{RGB}{120,120,120}
\definecolor{codestr}{RGB}{160,40,40}

\lstdefinestyle{shell}{
  language=bash,
  backgroundcolor=\color{codebg},
  basicstyle=\ttfamily\footnotesize,
  keywordstyle=\color{codekw}\bfseries,
  commentstyle=\color{codecm},
  stringstyle=\color{codestr},
  frame=single,
  breaklines=true,
  columns=fullflexible,
  keepspaces=true,
  showstringspaces=false,
  numbers=left,
  numberstyle=\tiny\color{codecm},
  xleftmargin=1.5em, xrightmargin=1em
}
\lstdefinestyle{perl}{
  language=Perl,
  backgroundcolor=\color{codebg},
  basicstyle=\ttfamily\footnotesize,
  keywordstyle=\color{codekw}\bfseries,
  commentstyle=\color{codecm},
  stringstyle=\color{codestr},
  frame=single,
  breaklines=true,
  columns=fullflexible,
  keepspaces=true,
  showstringspaces=false
}
\lstdefinestyle{smtp}{
  language={},
  backgroundcolor=\color{codebg},
  basicstyle=\ttfamily\footnotesize,
  frame=single,
  breaklines=true,
  columns=fullflexible,
  keepspaces=true,
  showstringspaces=true
}
\lstdefinestyle{yaml}{
  language=yaml,
  backgroundcolor=\color{codebg},
  basicstyle=\ttfamily\footnotesize,
  frame=single,
  breaklines=true,
  columns=fullflexible,
  keepspaces=true,
  showstringspaces=false
}
\lstdefinestyle{yargen}{
  language={},
  backgroundcolor=\color{codebg},
  basicstyle=\ttfamily\footnotesize,
  frame=single,
  breaklines=true,
  columns=fullflexible,
  keepspaces=true,
  showstringspaces=false
}
\lstdefinestyle{yar}{
  language={},
  backgroundcolor=\color{codebg},
  basicstyle=\ttfamily\footnotesize,
  frame=single,
  breaklines=true,
  columns=fullflexible,
  keepspaces=true,
  showstringspaces=false
}

\newcommand{\figdir}{../evidence/figures}

\begin{document}

\title{CVE-2026-73570: Full Lifecycle Evidence, Defender Detections, and
Remediation for Zimbra's SNMP-Sink Command Injection}

\author{Miguel Zabala\\[0.5ex]
\small Founder, Xpectra Security Research}

\markboth{CVE-2026-73570: Full Lifecycle Evidence, Defender Detections, and Remediation}
{Xpectra Security Research \MakeLowercase\expandafter\the\year}

\maketitle

\begin{abstract}
On 2026-06-26 the Zimbra security team published an advisory for an unauthenticated
remote code execution (RCE) in the SNMP monitoring component of Zimbra
Collaboration Suite (ZCS) versions before 10.1.20 (CVE-2026-73570, CVSS 3.1 8.9,
CWE-78). The vendor disclosure was intentionally limited: the trigger field on the
SMTP channel was not named. On 2026-08-17 CERT Polska (advisory 145/2026) reported
active exploitation in the wild, and CISA added the CVE to the KEV catalog on
2026-08-21 with a federal remediation deadline of 2026-08-24. This paper provides
what the public disclosure left out: a complete, reproducible, lab-verified
evidence chain from a crafted email to an unauthenticated shell, an empirical
identification of the trigger path, negative-control methodology, and a full
defender package — five Sigma rules, three YARA rules, three Suricata signatures,
EDR platform mappings, a non-invasive canary detector, an incident-response
playbook, and fix verification. All experiments were executed in an isolated,
zero-egress Docker lab on ZCS 10.1.0 GA, and the end-to-end pipeline was operated
by autonomous AI agents in a controlled environment, which is itself analyzed as a
case study in coherent long-horizon attack planning.
\end{abstract}

\begin{IEEEkeywords}
Cyberattacks, command injection, penetration testing, LLM agents,
incident detection, Zimbra, CVE-2026-73570.
\end{IEEEkeywords}

% =============================================================================
\section{Introduction}
\label{sec:intro}

Mail servers are trusted by the networks they sit in. Postfix logs every
envelope it processes — including recipient addresses it was never asked to
validate as data rather than text. When a monitoring subsystem consumes those
logs and interpolates captured fields into a shell command, the entire mail
pipeline becomes an unauthenticated command-execution channel. That is exactly
the shape of \textbf{CVE-2026-73570}, an OS command injection in Zimbra's
\texttt{zimbra-snmp} monitoring component that executes arbitrary commands as
the \texttt{zimbra} service user from a single crafted SMTP transaction — no
credentials, no open non-standard ports, and (on patched-default installs) no
SNMP listener reachable from the attacker at all.

The public record is strong but thin. The vendor advisory (2026-06-26) names the
component and severity; the fix (10.1.20, 2026-07-20) is a one-function rewrite
in the public \texttt{zm-build} repository (commit \texttt{c4837c66}); CERT
Polska (145/2026) documented active exploitation with an IOC of forged
``Service status change'' log lines; CISA KEV (2026-08-21) makes it a
priority-1 for federal systems. What no public source has provided is the
complete middle: the exact wire shape of the trigger, the process-level evidence
of execution, a methodology to verify the fix on a patched host, and detection
that defenders can deploy today.

\textbf{Our contributions.} Working in an isolated, zero-egress Docker lab with
ZCS 10.1.0 GA (a version inside the affected range), and operating the entire
pipeline with autonomous AI agents (Section~\ref{sec:agent}), we provide:
\begin{enumerate}
\item \textbf{An empirical trigger path.} The vendor did not name the SMTP
      field. We identify and prove the recipient-field path: a crafted quoted
      \texttt{RCPT TO:} local part survives postfix envelope handling verbatim,
      is captured by the \texttt{swatch} \texttt{watchfor} monitor, interpolated
      into the \texttt{doSNMP()} backtick command, and executed by
      \texttt{/bin/sh -c} as \texttt{zimbra} (Section~\ref{sec:exploit}).
\item \textbf{Full lifecycle evidence} from first packet to captured flag and
      defaced webmail, including a negative-control dry run that proves the sink
      executes code before any destructive payload is ever sent
      (Section~\ref{sec:exploit}).
\item \textbf{A trust-chain analysis} of why the injection is structurally
      inevitable on unpatched versions — which hop breaks the
      log-as-data assumption and what the vendor patch changes at the
      \texttt{execve} level (Section~\ref{sec:vuln}).
\item \textbf{A complete defender package}: 5 Sigma rules, 3 YARA rules, 3
      Suricata signatures, EDR mappings (Microsoft Defender for Linux,
      CrowdStrike), a non-invasive canary detector, an IOC set, a remediation
      guide, and an incident-response playbook — with each rule demonstrated
      firing against our own exploitation (Section~\ref{sec:detection}).
\item \textbf{Fix verification} by re-applying the vendor patch to a second lab
      instance and showing the identical exploit fails while a positive control
      on the stock instance still fires (Section~\ref{sec:fix}).
\item \textbf{A case study} in autonomous attack planning: a beat-logged,
      think-then-act agent pipeline with a live operations map, analyzed against
      the long-horizon-coherence limitations documented in recent LLM
      pentesting research~\cite{wang2025checkmate} (Section~\ref{sec:agent}).
\end{enumerate}

The remainder of the paper: Section~\ref{sec:vuln} characterizes the
vulnerability and the trust chain; Section~\ref{sec:lab} describes the lab and
methodology; Section~\ref{sec:exploit} presents the exploitation evidence;
Section~\ref{sec:fix} the fix verification; Section~\ref{sec:detection} the
defender package; Section~\ref{sec:remediation} remediation and IR; and
Section~\ref{sec:ethics} ethics, reproducibility, and the agent-pipeline
discussion. All artifacts (exploit, detector, rules, lab) are released in the
companion folder \texttt{CVE-2026-73570/} of our Research repository.

% =============================================================================
\section{Vulnerability Characterization}
\label{sec:vuln}

\subsection{The Zimbra SNMP monitoring stack}

ZCS ships an optional SNMP monitoring package (\texttt{zimbra-snmp}) whose
purpose is to emit SNMP traps when Zimbra log conditions change. The moving
parts, all running on the Zimbra host as the unprivileged \texttt{zimbra}
user (uid 1000):

\begin{itemize}
\item \textbf{swatch/swatchdog}: a Perl log-monitoring tool. \texttt{swatchdog}
taps log files; each \texttt{watchfor} pattern in
\texttt{/opt/zimbra/conf/swatchrc} can capture regex groups, and an
associated \texttt{donotify} action is invoked with the captured values.
\item \textbf{doSNMP()}: a Perl function defined in \texttt{swatchrc}
(\texttt{perlcode}) that the notify actions call. It assembles and runs a
\texttt{snmptrap} command carrying the captured \texttt{SERVICE} value.
\item \textbf{The log pipeline}: Postfix delivery records (including
\texttt{to=<recipient>} lines) are copied by rsyslog into
\texttt{/var/log/zimbra.log}, which \texttt{swatchdog} tails in real time.
\end{itemize}

The enabling configuration: \texttt{zmlocalconfig snmp\_notify=1} with a
trap destination. The precondition is documented by CERT
Polska~\cite{certpolska} and by the vendor advisory; our lab runs it exactly
that way (Section~\ref{sec:lab}).

\subsection{The vulnerable code}

In unpatched ZCS (we use 10.1.0 GA, build 4655), the installed
\texttt{swatchrc} contains:

\begin{lstlisting}[style=perl,caption={Vulnerable doSNMP() (ZCS 10.1.0, /opt/zimbra/conf/swatchrc)}]
perlcode 0
sub dosnmp {
    my %args = (@_);
    print "SNMP notification: $args{MESSAGE}\n";
    `$snmptrap $snmpsvctrap $snmpsvcname
      s $args{SERVICE}
      $snmpsvcstatus i $statuses{$args{STATUS}}`;
}
\end{lstlisting}

The command is executed inside \textbf{Perl backticks}. Perl expands
\texttt{\$args\{SERVICE\}} into the command string \emph{first}, and the
resulting string is then handed to \texttt{/bin/sh -c} for execution. Any
shell metasyntax present in the value of \texttt{SERVICE} is therefore
executed: command substitution \texttt{\$(...)}, pipes, redirects, chained
commands.

The vendor fix (10.1.20, commit \texttt{c4837c66}, public) rewrites the same
function:

\begin{lstlisting}[style=perl,caption={Fixed doSNMP() (10.1.20, Zimbra/zm-build c4837c66)}]
perlcode 0
sub dosnmp {
    my %args = (@_);
    print "SNMP notification: $args{MESSAGE}\n";
    system("/opt/zimbra/common/bin/snmptrap", "-v", "2c",
           "-c", "zimbra", $traphost, "",
           $snmpsvctrap, $snmpsvcname, "s",
           $args{SERVICE}, $snmpsvcstatus,
           "i", $statuses{$args{STATUS}});
}
\end{lstlisting}

This is the multi-argument form of \texttt{system()}, which calls
\texttt{execve()} directly: there is no shell, no metacharacter parsing, and
the attacker-controlled value becomes one inert \texttt{argv} element passed
to \texttt{snmptrap}. Note the subtle operational detail our patch
re-application surfaced: \texttt{zmswatchctl} re-renders the active
\texttt{swatchrc} \emph{from} \texttt{swatchrc.in} on every SNMP service
restart, so a complete fix (or a complete lab revert) must touch both files.

\subsection{The trust chain, hop by hop}

\begin{figure*}[t]
\centering
\includegraphics[width=\textwidth]{\figdir/fig13-attack-chain.png}
\caption{The complete chain from crafted SMTP to \texttt{/bin/sh -c}. No
hop between the wire and the shell sanitizes the recipient text. The red box
shows the vulnerable sink; the green box shows the 10.1.20 fix (multi-arg
\texttt{system()} = \texttt{execve}, no shell).}
\label{fig:chain}
\end{figure*}

\textbf{Hop 1 — the wire.} The attacker sends a normal SMTP session to port
25 from an address the MTA will accept (Section~\ref{sec:precond}). The
recipient's local part carries the payload inside a quoted address
(Section~\ref{sec:exploit}). Nothing on the wire looks like a command — it is
an e-mail address.

\textbf{Hop 2 — the MTA logs verbatim.} Postfix writes its delivery decision,
including the full \texttt{to=<...>} address, into the mail log; rsyslog
copies it into \texttt{/var/log/zimbra.log}. Postfix treats the address as an
opaque token (with the documented restrictions on \texttt{<}, \texttt{>}, and
spaces — Section~\ref{sec:exploit}); it does not need to be a deliverable
address at all. \emph{This is the first trust-break: the recipient field is
attacker text, but every downstream consumer treats it as benign data.}

\textbf{Hop 3 — the monitor captures it.} A \texttt{watchfor} pattern in
\texttt{swatchrc} matches the log line and captures the address into
\texttt{\$1}. The monitor's job is pattern matching on untrusted text —
appropriate for its purpose — but the \emph{action} attached to that capture
decides the security outcome.

\textbf{Hop 4 — the action interpolates.} \texttt{donotify SERVICE=\$1}
hands the captured address to \texttt{doSNMP()}.

\textbf{Hop 5 — the sink executes.} The backtick command is built by Perl
string interpolation and executed by \texttt{/bin/sh -c}. The
attacker-controlled \texttt{\$(...)} now runs. \emph{This is the trust-break
the vendor patched}: the fix moves the value from ``part of a shell command
string'' to ``one argv element'' — an \texttt{execve} boundary where
metasyntax has no meaning.

\textbf{Consequence.} Execution is as \texttt{zimbra}: full read access to
all mailboxes (PII/ATO potential), the ability to inject mail from the
domain, a pivot point on a typically network-privileged host, and trivial
visible defacement (we demonstrate swapping the webmail/admin login logos).
Privilege escalation beyond \texttt{zimbra} is out of scope here; the KEV
listing reflects the impact at the \texttt{zimbra} level alone.

\subsection{Preconditions and the AC:H}
\label{sec:precond}

The CVSS vector \texttt{AV:N/\hspace{0pt}AC:H/\hspace{0pt}PR:N/\hspace{0pt}UI:N/\hspace{0pt}S:C/\hspace{0pt}C:H/\hspace{0pt}I:H/\hspace{0pt}A:L} marks
Attack Complexity \texttt{High}. Our lab analysis attributes that to the
reachable-SMTP precondition: the crafted address must reach the MTA.
Realistic paths are (i) any address inside Postfix \texttt{mynetworks};
(ii) an authenticated mail-submission account — any mailbox, since
\texttt{PR:N} means no privilege and a low-value account suffices; or
(iii) an internal relay. In our lab, path~(i) is instantiated by the
internal bridge — a documented lab condition that mirrors how such sinks
are reached from the inside in real deployments.
The complexity is in reaching the MTA with an address that is logged and
then bounced (no deliverable mailbox is required). Once inside that path,
no user interaction, no authentication to any privileged service, and no
non-standard port is needed — SMTP 25 is the door.

The remaining preconditions (all default- or near-default in real
deployments that use SNMP monitoring) are: \texttt{zimbra-snmp} installed,
\texttt{snmp\_notify=1}, \texttt{swatchdog} running against
\texttt{/var/log/zimbra.log}, and a \texttt{watchfor} whose action feeds the
captured text into \texttt{doSNMP}. Section~\ref{sec:lab} documents our lab
instantiation of each.
